\section{Truncation transport and canonical normalization}
\label{rev:truncations}
\label{sec:general-pullbacks}
The analytic transport is \eqref{eq:general-pullback-transport}; its
gauge and logarithmic derivatives are denoted here by $g,u,v$.
All applications below retain the nondegeneracy and branch conditions
of the general theory. We now prove the additional finite-polynomial
condition needed to transport raw weighted truncations.
\subsection{A polynomial criterion for truncation transport}

Write $[H]_{<p}=\sum_{n=0}^{p-1}[t^n]H\,t^n$.
All congruences in the next lemma are coefficientwise, followed by
reduction to $\mathbb F_p[[t]]$ or $\mathbb F_p(t)$ as appropriate.

\begin{lemma}[Finite Frobenius transformation]
\label{lem:finite-frobenius-pullback}
Let $p>3$, let $g\in1+t\mathbb Z_p[[t]]$, let
$r\in t\mathbb Z_p[[t]]$, and suppose
$H=gF_s(r)\in\mathbb Z_p[[t]]$.
If
\[
 \Lambda_p(t)=g(t)^{1-p}T_p(r(t))
\]
reduces to a polynomial of degree less than $p$, then
\begin{equation}\label{eq:finite-frobenius-pullback}
 [H]_{<p}(t)\equiv\Lambda_p(t)\pmod p.
\end{equation}
If $u=\delta g/g$ and $v=\delta r/r$, then
\begin{equation}\label{eq:finite-weighted-pullback}
 (A+B\delta)[H]_{<p}(t)
 \equiv g(t)^{1-p}
 \bigl\{(A+Bu(t))T_p(r(t))
        +Bv(t)r(t)T'_p(r(t))\bigr\}\pmod p.
\end{equation}
The last identity is valid wherever its reduced rational expressions
are defined.
\end{lemma}
\begin{proof}
The rational parameters $1/2,s,1-s$ are $p$-adic integers, so their
binomial coefficients show that all coefficients of $F_s$ are
$p$-integral.  Since $r$ has positive order,
$F_s(r)-T_p(r)$ is divisible by $t^p$.  In characteristic $p$,
\[
 g(t)^{-p}=g(t^p)^{-1}=1+O(t^p).
\]
It follows that $g^{1-p}T_p(r)$ and $gF_s(r)$ have the same
coefficients of degree less than $p$.  The assumed polynomial degree
bound proves \eqref{eq:finite-frobenius-pullback}.  Differentiating
in characteristic $p$ gives
\[
 \frac{\delta(g^{1-p})}{g^{1-p}}=(1-p)u=u,
\]
and then the chain rule proves
\eqref{eq:finite-weighted-pullback}.
\end{proof}

\begin{proposition}[Three exact finite transformations]
\label{prop:three-finite-transformations}
For every prime $p>3$, the following are polynomial congruences of
degree less than $p$:
\begin{align}
 [\mathcal G]_{<p}(t)
 &\equiv(1+Kt)^{p-1}
       T_p\!\left(\frac{Rt}{(1+Kt)^2}\right),
       &&R=4K,\label{eq:finite-quadratic}\\
 [H_{\mathrm{Domb}}]_{<p}(t)
 &\equiv(1-4t)^{p-1}
       T_{p,1/3}\!\left(\frac{108t^2}{(1-4t)^3}\right),
       \label{eq:finite-domb}\\
 [H_9]_{<p}(t)
 &\equiv D(t)^{(p-1)/2}
       T_{p,1/6}\!\left(\frac{1728t^3}{D(t)^3}\right),
       &&D(t)=1-6t-27t^2.\label{eq:finite-nine}
\end{align}
Here the first line uses the corresponding value of $s$, and
$T_{p,s}$ denotes the truncation of $F_s$.
Their weighted forms are obtained by
\eqref{eq:finite-weighted-pullback}, with exactly the logarithmic
derivatives in the analytic transformations.
\end{proposition}
\begin{proof}
In the reduced truncations, a Pochhammer numerator first vanishes
after degrees
\[
 d_s(p)=
 \begin{cases}
 \lfloor p/6\rfloor,&s=1/6,\\
 \lfloor p/4\rfloor,&s=1/4,\\
 \lfloor p/3\rfloor,&s=1/3,\\
 (p-1)/2,&s=1/2.
 \end{cases}
\]
Thus $\deg T_{p,s}=d_s(p)$.  In
\eqref{eq:finite-quadratic}, the $n$th term is a constant times
$t^n(1+Kt)^{p-1-2n}$, whose exponent is nonnegative and whose degree
is at most $p-1$.  In \eqref{eq:finite-domb}, its counterpart is
$t^{2n}(1-4t)^{p-1-3n}$, again with nonnegative exponent and degree
at most $p-1$.  In \eqref{eq:finite-nine}, it is
$t^{3n}D^{(p-1)/2-3n}$; the exponent is nonnegative because
$n\le\lfloor p/6\rfloor$, and its degree is $p-1-3n$.
Each right side therefore satisfies the polynomial degree condition
of Lemma~\ref{lem:finite-frobenius-pullback}.  The defining analytic
identities and that lemma prove all three congruences.
\end{proof}

The arithmetic prefactors are essential.  For example, the factor
in \eqref{eq:finite-domb} is $(1-4t)^{p-1}$, whereas the analytic
factor is $(1-4t)^{-1}$.  Their logarithmic derivatives agree modulo
$p$, which is precisely what permits linear weights to transfer.

\begin{theorem}[Arithmetic classification for the additional families]
\label{thm:additional-raw-density}
In each of the quadratic-companion, Domb, and level-nine families,
consider integer denominators and primitive integer linear weights
for which the defining series converges ordinarily.  Let
\[
 Z_{A,B,t_0}=
 \left\{p>3:(A+B\delta)[H]_{<p}(t_0)\equiv0\pmod p\right\},
 \qquad t_0=1/C,
\]
omitting the finitely many primes at which the data are not integral
or the displayed rational transformations have poles.  The following
are equivalent:
\begin{enumerate}
\item The lower natural density of $Z_{A,B,t_0}$ is greater than $3/4$.
\item The set $Z_{A,B,t_0}$ contains all but finitely many primes.
\item Up to simultaneous sign of the weights, the data are one of the
five quadratic-companion rows, the two Domb rows, or the two level-nine
rows in the respective CM classification.
\end{enumerate}
\end{theorem}
\begin{proof}
At a nonendpoint convergent parameter, the analytic transformations
have $g\ne0$, $v\ne0$, and rational $r,u,v$.
Proposition~\ref{prop:three-finite-transformations} identifies the
raw weighted truncation, up to a nonzero finite-field factor, with
the standard weighted truncation having the fixed rational weights
\[
 (a,b)=(A+Bu(t_0),Bv(t_0)).
\]
The sets of vanishing rational primes therefore differ by at most
a finite set.  Theorem~\ref{thm:weighted-density-cm} shows that the
first condition forces a CM fiber.  The exact CM parameter
classification leaves the denominators already listed.
Proposition~\ref{prop:cm-weighted-slope} then identifies the unique
slope of density greater than $3/4$ with the CM eigenline.  Its period
determinant characterization and invertible analytic transport
identify this slope with the displayed complex identity.  At that
slope the congruence holds outside a finite set of primes, proving
the second condition.  Conversely, every displayed identity selects
that line and therefore satisfies the second condition.  The second
condition plainly implies the first.

There are two kinds of convergent endpoint not covered by this
regular argument: $C=K$ for a quadratic companion, and $C=16$ for
Domb.  In either case convergence forces $B=0$, and the transformed
parameter is $z=1$.  The finite transformations reduce the zero
condition to $T_p(1)=0$.  Finite Clausen gives
$T_p(1)=U_p(1/2)^2$, and the Hasse calculation remains valid at
$x=1/2$, where the elliptic fiber is smooth and CM.  Deuring's
criterion therefore gives density $1/2$ for these zero sets.  They
do not satisfy the first condition and supply no additional rows.
\end{proof}

This theorem is unconditional, but its hypothesis is arithmetic.
It does not derive the weighted congruences from a complex identity
at an unknown non-CM fiber.  Under the appropriately enlarged
quadratic-period hypothesis, the analytic classifications and these
arithmetic classifications coincide.


\subsection{Why the regularity conditions cannot be suppressed}
\label{sec:countermodel-guardrails}
The next examples concern rational and algebraic transforms of the
standard functions. They do not contradict the standard-branch
classification. They show precisely why a general formulation must
control zeros of the gauge, critical points of the parameter map, the
coefficient field, and singular source parameters.

Suppose $H(t)=g(t)F(\phi(t))$ and put $\delta=t\,d/dt$.
At an algebraic point $t_0$ where $F$ is analytic at $\phi(t_0)$,
$g,\phi$ are regular, and $t_0g(t_0)\phi(t_0)\ne0$, the transport is
\begin{equation}\label{eq:countermodel-transport}
 (A+B\delta)H=g\left[
 \left(A+B\frac{\delta g}{g}\right)F
 +B\frac{\delta\phi}{\phi}\thetaop F\right].
\end{equation}
Its determinant is $g(t_0)^2t_0\phi'(t_0)/\phi(t_0)$.
Thus nonvanishing of this determinant transfers algebraic projective
uniqueness. Rational projective slopes require a further field check.

\begin{proposition}[Critical pullbacks with cubic supercongruences]
\label{prop:critical-countermodels}
Let $m>3$ be an integer with $m\nmid64$, and write
\[
 H_m(t)=F_{1/2}\!\left(\frac{2t-t^2}{m}\right)
       =\sum_{k\ge0}h_kt^k.
\]
The series and its differentiated series converge absolutely at
$t=1$. The elliptic fiber there is smooth and non-CM, but
\[
 \delta H_m(1)=0,\qquad
 \sum_{k=0}^{p-1}kh_k\equiv0\pmod{p^3}
 \quad\text{for every odd prime }p\nmid m.
\]
\end{proposition}
\begin{proof}
Put $\phi_m(t)=(2t-t^2)/m$. Choose $r>1$ with $2r+r^2<m$.
Then $|\phi_m(t)|<1$ on $|t|\le r$, proving convergence.
The derivative vanishes because $\phi'_m(1)=0$ while $F_{1/2}$ is
analytic at $1/m$. The invariant is
\[
 j(1/m)=\frac{64(4m-1)^3}{m}.
\]
Since $\gcd(m,4m-1)=1$ and $m\nmid64$, it is a nonintegral rational
number. It cannot be a CM invariant, because singular moduli are
algebraic integers.

For the congruence write
$c_n=\binom{2n}{n}^3/64^n$ and
\[
 V_p(t)=\sum_{n=0}^{p-1}c_n\phi_m(t)^n.
\]
The coefficients of degree less than $p$ agree with those of $H_m$.
Terms with $n\le(p-1)/2$ have degree at most $p-1$, while terms with
$(p+1)/2\le n<p$ have $p^3\mid c_n$ in $\Z_{(p)}$. Hence
\[
 V_p(t)-\sum_{k<p}h_kt^k\in p^3\Z_{(p)}[t].
\]
Since $V'_p(1)=0$ exactly, differentiation and evaluation at 1 give
the asserted congruence.
\end{proof}

This example has the primitive input pair $(A,B)=(0,1)$ and the
allowed multiplier $\alpha=0$. It does not refute a CM statement
restricted to nonzero multipliers. It does refute an unrestricted
claim that cofinite weighted supercongruences force CM for every
rational pullback of an elliptic period square. At this critical
point the proposed de Rham direction is zero, so there is no
complementary line to which the algebraization theorem could apply.
The condition $B\ne0$ cannot replace that geometric requirement.
Multiplication by $1+t$ gives the non-CM zero identity with primitive
pair $(-1,2)$; this last observation concerns the complex value only.

Vanishing gauges can destroy even uniqueness. The normalized rational
coefficient function
\[
 J(t)=(1-t)^2F_{1/2}(t/5)
\]
has radius at least 5 and satisfies $J(1)=\delta J(1)=0$, at the
smooth non-CM fiber $j=438976/5$. Every primitive integer pair gives
zero there, so one parameter has infinitely many slopes.
Separately, a smoothness condition cannot be replaced by a
noncriticality condition: for $m>2$,
\[
 K_m(t)=F_{1/2}\!\left(\frac{t(1-t)}m\right),\qquad
 (1+8m\delta)K_m(1)=0.
\]
Here the pullback is noncritical at 1, but its image is the nodal
parameter $z=0$; indeed $K_m(1)=1$ and $\delta K_m(1)=-1/(8m)$.

\begin{proposition}[Rational Taylor coefficients do not preserve rational slopes]
\label{prop:gauge-rational-slope-countermodel}
The rational-coefficient function
\[
 L(t)=\frac{1+\sqrt{1+t/3}}2\,F_{1/2}(t/4),
\]
with square root equal to 1 at zero, is normalized by $L(0)=1$ and
converges absolutely, with its derivative, at $t=1$. Its fiber there
has CM. Nevertheless it has no nonzero rational pair $(A,B)$ for an
inverse-$\pi$ identity at $t=1$.
\end{proposition}
\begin{proof}
The scalar $g(t)=(1+\sqrt{1+t/3})/2$ has rational Taylor coefficients
and is analytic for $|t|<3$. The other factor is analytic for $|t|<4$.
The standard identity at $z=1/4$ is
$(1+6\thetaop)F_{1/2}(1/4)=4/\pi$, and its invariant is $j=54000$.
By algebraic projective uniqueness the unique slope is $1/6$.
But $\delta g(1)/g(1)=(2-\sqrt3)/4$, so
\eqref{eq:countermodel-transport} transports that slope to
\[
 \frac AB=\frac16-\frac{2-\sqrt3}{4}
 =-\frac13+\frac{\sqrt3}{4}\notin\Q.
\]
Taking $B=1$ gives multiplier $2g(1)/3\ne0$. Every nonzero pair is
on this unique algebraic slope, so none is rational.
\end{proof}

A critical pullback over a singular source parameter requires a
different local calculation. In particular it need not give a zero
derivative. Clausen's transformation gives the germ identity
\[
 M(t)=F_{1/2}(t(2-t))
 ={}_2F_1(1/2,1/2;1;t/2)^2.
\]
The right side has radius 2. Consequently the series and derivative
converge absolutely at $t=1$. In complete elliptic-integral parameter
notation, $f(x)={}_2F_1(1/2,1/2;1;x)=2K(x)/\pi$, and at $x=1/2$
\[
 f'(1/2)=\frac2\pi(2E(1/2)-K(1/2)),\qquad
 2K(1/2)E(1/2)-K(1/2)^2=\frac\pi2.
\]
It follows that
\[
 \delta M(1)=f(1/2)f'(1/2)=\frac2\pi.
\]
The fiber has $j=1728$ and is CM. The source derivative $F'_{1/2}(z)$
is singular at $z=1$, so direct substitution into a product of the
vanishing pullback derivative with $F'_{1/2}(1)$ was never legitimate.
This example also shows that an algebraic pullback can improve the
ordinary convergence of a differentiated series.

\subsection{The different behavior of raw Euler truncations}
Finite transformation is not an automatic consequence of
analytic transformation.  For the linear Euler companion
$H=F_s/(1-t)$, direct summation of coefficients gives instead
\[
 [H]_{<p}(t)
  =\frac{T_p(t)-t^pT_p(1)}{1-t}\pmod p.
\]
The extra boundary term is generally nonzero.  In this case
$(1-t)^{p-1}T_p(t)$ need not have degree less than $p$, so
Lemma~\ref{lem:finite-frobenius-pullback} does not apply.  Transported
standard truncations still provide an arithmetic criterion for its
35 slopes, but identifying that criterion with the raw Euler
truncations is false in general, as the next exact example shows.

\begin{proposition}[A CM Euler row with raw zero density $1/2$]
\label{prop:euler-raw-density-obstruction}
For $s=1/2$, $t_0=1/4$, and $H=F_{1/2}/(1-t)$, the genuine CM
identity
\[
 (-1+6\delta)H(1/4)=\frac{16}{3\pi}
\]
has raw weighted-truncation zeros of natural density exactly $1/2$.
\end{proposition}
\begin{proof}
The exact formula for $[H]_{<p}$, differentiated in characteristic
$p$, gives
\[
 (-1+6\delta)[H]_{<p}(1/4)
 \equiv\frac43(1+6\thetaop)T_p(1/4)-\frac13T_p(1).
\]
The first term vanishes at all but finitely many primes by the CM
line criterion for the standard row $(1,6,256,4)$.
Put $m=(p-1)/2$.  Finite Clausen and the hypergeometric expression
for the Legendre polynomial give
\[
 T_p(1)=U_p(1/2)^2=P_m(0)^2\pmod p.
\]
For odd $m$, $P_m(0)=0$; for even $m$,
\[
 P_m(0)=(-1)^{m/2}2^{-m}\binom{m}{m/2}\not\equiv0\pmod p.
\]
Thus, outside finitely many primes, the raw Euler weighted
truncation vanishes exactly when $p\equiv3\pmod4$.  These primes
have natural density $1/2$.
\end{proof}


\subsection{Finite jets do not determine raw congruence densities}
\label{subsec:C-gauge-obstruction}
The arithmetic transfer problem cannot be solved by retaining only
global boundedness, irreducible elliptic symmetric-square monodromy,
and finitely many derivatives at the identity point.  The following
construction preserves all these data but changes the prime-zero
density from one to zero.

\begin{proposition}\label{prop:C-finite-jet-density-zero}
Put $F=F_{1/2}$ and, for any integer $M\ge2$, define
\[
 g_M(z)=1+\frac{z(4z-1)^M}{3-4z},\qquad H_M(z)=g_M(z)F(z).
\]
Then $H_M(0)=1$, $H_M(192t)\in\mathbb Z[[t]]$, and the minimal
differential module of $H_M$ is isomorphic over $\mathbb Q(z)$ to
the irreducible rank-three differential module of $F$.
All derivatives of $H_M$ and $F$ through order $M-1$ agree at
$z_0=1/4$.  In particular, an ordinarily absolutely convergent
identity holds:
\[
 (1+6\thetaop)H_M(1/4)=\frac4\pi.
\]
Nevertheless, if $H_{M,p}$ is the Taylor polynomial of $H_M$ of
degree less than $p$, the set of primes for which
\[
 (1+6\thetaop)H_{M,p}(1/4)\equiv0\pmod p
\]
has natural density zero.
\end{proposition}
\begin{proof}
Write
\[
 T_p(z)=\sum_{n=0}^{p-1}\frac{\binom{2n}{n}^3}{64^n}z^n
 \quad\text{in }\mathbb F_p[z],\qquad
 P_M(z)=3-4z+z(4z-1)^M.
\]
For $p>3$, the polynomial $T_p$ has degree at most $(p-1)/2$.
If $p\ge2M+3$, then $\deg(P_MT_p)\le p-1$.
Finite geometric summation therefore gives the polynomial identity
\begin{equation}\label{eq:C-gauge-finite}
 H_{M,p}(z)=
 \frac{P_M(z)T_p(z)-(4z/3)^pP_M(3/4)T_p(3/4)}{3-4z}
 \qquad\text{in }\mathbb F_p[z].
\end{equation}
Indeed the numerator vanishes at $z=3/4$, the quotient has degree
at most $p-1$, and its germ agrees modulo $z^p$ with
$P_M(z)F(z)/(3-4z)$.

Now $g_M(z_0)=1$ and $g'_M(z_0)=0$.  Differentiation in
characteristic $p$ annihilates $(4z/3)^p$; also
\[
 \left.(1+6\thetaop)\frac1{3-4z}\right|_{z=z_0}=2,
 \qquad (4z_0/3)^p=\frac13,
 \qquad P_M(3/4)=\frac34\,2^M.
\]
Consequently \eqref{eq:C-gauge-finite} implies
\begin{equation}\label{eq:C-gauge-weighted}
 (1+6\thetaop)H_{M,p}(1/4)
 =(1+6\thetaop)T_p(1/4)-2^{M-1}T_p(3/4)
 \quad\text{in }\mathbb F_p.
\end{equation}
The standard CM row $(A,B,C)=(1,6,256)$ and its proved CM
congruence give $(1+6\thetaop)T_p(1/4)=0$ outside a finite set
of primes.  Finite Clausen, applied with $x=1/4$ and
$4x(1-x)=3/4$, identifies $T_p(3/4)$ with the square of
the Hasse polynomial of the Legendre curve
\[
 E:\quad y^2=x(x-1)(x-1/4).
\]
Thus $T_p(3/4)=0$ precisely at the supersingular good primes.
The invariant
\[
 j(E)=256\frac{(1-1/4+1/16)^3}{(1/16)(9/16)}
     =\frac{35152}{9}
\]
is not an algebraic integer; hence $E$ is non-CM.  Its
supersingular rational primes have density zero by the non-CM
case of Theorem~\ref{thm:B-qcurve-density} (or directly by Serre's
open-image theorem and Chebotarev).  Equation
\eqref{eq:C-gauge-weighted} proves the density assertion.

Finally,
\[
 F(192t)=\sum_{n\ge0}\binom{2n}{n}^3 3^n t^n,
 \qquad
 g_M(192t)=1+\frac{64t(768t-1)^M}{1-256t}
\]
have integral coefficients.  Rational multiplication by the nonzero
function $g_M$ is an invertible change of differential module;
it preserves rank and irreducibility over $\mathbb Q(z)$.
The difference $H_M-F$ vanishes to order $M$ at $z_0$.
Moreover, $H_M$ is analytic for $|z|<3/4$, so evaluation and
differentiation at $z_0$ are ordinary absolutely convergent
operations.  This proves every claimed property.
\end{proof}

\begin{corollary}\label{cor:C-two-jet-density-zero}
The same obstruction persists when arbitrarily many Taylor
coefficients at the cusp are prescribed as well.  For $M\ge2$ and
$N\ge1$, put
\[
 g_{M,N}(z)=1+\frac{z^N(4z-1)^M}{3-4z},\qquad
 H_{M,N}=g_{M,N}F.
\]
Then $H_{M,N}-F$ vanishes to order $N$ at $0$ and to order $M$
at $1/4$.  All conclusions of
Proposition~\ref{prop:C-finite-jet-density-zero} hold, with the
finite formula for every $p\ge2(M+N)+1$ replaced by
\[
 (1+6\thetaop)H_{M,N,p}(1/4)
 =(1+6\thetaop)T_p(1/4)
   -2^{M+1-2N}3^{N-1}T_p(3/4)\pmod p.
\]
\end{corollary}
\begin{proof}
Replace $P_M$ in the preceding proof by
$3-4z+z^N(4z-1)^M$, of degree $M+N$.
Its value at $3/4$ is $(3/4)^N2^M$, giving the displayed
coefficient.  It is nonzero modulo every $p>3$.
The rescaling $z=192t$ again gives integral coefficients, since
$192^N/3$ is integral.  The other assertions follow from the
indicated orders of vanishing and the preceding density proof.
\end{proof}

This proposition changes the generating function and is not a
counterexample to CM forcing for any of the four standard branches.
It shows that a successful transfer theorem must retain the specified
Frobenius-compatible coefficient normalization; no prescribed finite
jet at the complex identity point can replace that normalization.

\subsection{A normalization retained by the standard coefficients}
\label{subsec:C-canonical-lucas}

The preceding rational gauges preserve finite complex data but lose
an exact arithmetic property of the standard coefficients: Lucas
congruences. This supplies a restriction on possible transfer
arguments, although it does not establish a transfer theorem.

\begin{proposition}[Standard Lucas normalization]
\label{prop:C-standard-lucas}
Let $u_\ell(n)$ be the factorial coefficients in \eqref{eq:series},
and put $k_\ell=6,4,3,2$ for $\ell=1,2,3,4$, respectively.
For every prime $p$, $a\geq0$, and $0\leq b<p$,
\[
 u_\ell(pa+b)\equiv u_\ell(a)u_\ell(b)\pmod p.
\]
If $p>k_\ell$, then
\[
 T_{\ell,p}(t)=\sum_{n=0}^{p-1}u_\ell(n)t^n\in\mathbb F_p[t]
 \quad\text{has degree}\quad
 d_\ell(p)=\left\lfloor\frac{p-1}{k_\ell}\right\rfloor,
 \qquad 1\leq d_\ell(p)<p-1.
\]
Writing $G_\ell(t)=\sum_{n\geq0}u_\ell(n)t^n$, one has
\begin{equation}\label{eq:C-standard-lucas-functional}
 G_\ell(t)=T_{\ell,p}(t)G_\ell(t^p)
 \qquad\text{in }\mathbb F_p[[t]].
\end{equation}
The same conclusions hold for $F_{s_\ell}(z)=G_\ell(z/R_\ell)$
when $p\nmid R_\ell$.
\end{proposition}
\begin{proof}
Legendre's factorial valuation formula gives
\[
 v_p(u_\ell(n))=\sum_{j\geq1}\Delta_\ell(n/p^j),
\]
where the $\Delta_\ell$ are periodic of period one and on $[0,1)$ equal
\[
 \lfloor6x\rfloor-\lfloor3x\rfloor,\quad
 \lfloor4x\rfloor,\quad
 \lfloor3x\rfloor+\lfloor2x\rfloor,\quad
 3\lfloor2x\rfloor.
\]
They are nonnegative everywhere and positive on $[1/k_\ell,1)$.
If $b\geq p/k_\ell$, the first summand shows that both
$u_\ell(pa+b)$ and $u_\ell(b)$ vanish modulo $p$.
If $b<p/k_\ell$, there are no low-digit carries in the respective
integer presentations
\[
 \binom{6n}{3n}\binom{3n}{n,n,n},\qquad
 \binom{4n}{n,n,n,n},\qquad
 \binom{3n}{n,n,n}\binom{2n}{n},\qquad
 \binom{2n}{n}^{\!3}.
\]
Lucas's multinomial theorem proves the required congruence.
For $0\leq n<p/k_\ell$, every factorial argument is below $p$,
so $u_\ell(n)$ is a unit; for $p/k_\ell\leq n<p$ it is divisible
by $p$. This proves the degree formula. Coefficientwise digit
decomposition gives \eqref{eq:C-standard-lucas-functional}.
Rescaling by $R_\ell$ preserves the congruence by Fermat's theorem.
\end{proof}

\begin{proposition}[Rigidity under rational scalar gauges]
\label{prop:C-lucas-rational-rigidity}
Let $f\in\mathbb F_p[[t]]$ satisfy $f(0)=1$ and
\[
 f(t)=T(t)f(t^p),\qquad T\in\mathbb F_p[t],\quad
 1\leq\deg T<p-1.
\]
Let $g\in\mathbb F_p(t)$ be regular at zero with $g(0)=1$.
If $h=gf$ also satisfies Lucas congruences, then $g=1$.
Consequently, any nonconstant rational $g\in\Q(t)$ with $g(0)=1$
destroys the Lucas property of each standard series at every
sufficiently large prime.
\end{proposition}
\begin{proof}
Write $g=P/Q$ with $P,Q$ coprime and $P(0)=Q(0)\ne0$.
Lucas congruences for $h$ give $h(t)=U(t)h(t^p)$ with
$U\in\mathbb F_p[t]$ and $\deg U<p$. Since $g(t^p)=g(t)^p$,
\[
 U=Tg^{1-p}=\frac{TQ^{p-1}}{P^{p-1}}.
\]
Thus $P^{p-1}$ divides $T$, and $\deg T<p-1$ makes $P$ constant.
If $Q$ were nonconstant, then
$\deg U=\deg T+(p-1)\deg Q\geq p$, a contradiction.
Normalization gives $g=1$.
For a rational gauge in characteristic zero, exclude the finitely
many primes at which its coefficients are not integral, its numerator
or denominator degree drops, their resultant vanishes, or the
standard truncation has degree zero. Its reduction remains
nonconstant at every other prime, so the first assertion applies.
\end{proof}

There is a corresponding correction to raw truncation. For a normalized
series $h$ in characteristic $p$, define its formal Frobenius quotient by
\[
 \mathcal Q_p(h)(t)=\frac{h(t)}{h(t^p)}.
\]
In the notation of Proposition~\ref{prop:C-lucas-rational-rigidity},
\begin{equation}\label{eq:C-lucas-gauge-covariance}
 \mathcal Q_p(gf)=g^{1-p}T.
\end{equation}
The right side is rational, and its Taylor polynomial of degree less
than $p$ is the raw truncation of $gf$, since $h(t^p)\equiv1\pmod{t^p}$.
If $g(t_0)=1$ and $g'(t_0)=0$ at a regular point $t_0\in\mathbb F_p$,
then the value and first derivative of this rational quotient equal
those of $T$ at $t_0$. Thus
\[
 (A+Bt\partial_t)\mathcal Q_p(gf)(t_0)
 =(A+Bt\partial_t)T(t_0).
\]
This explains exactly where the gauge examples change the raw
weighted congruence: a rational Frobenius tail is discarded.
It is an identity in characteristic $p$, not a comparison theorem
for the original complex value.

\begin{remark}[Algebraic gauges remain possible]
\label{rem:C-algebraic-lucas-gauge}
The rationality hypothesis in
Proposition~\ref{prop:C-lucas-rational-rigidity} is necessary.
For $g(z)=(1-z)^{-1/2}$ and every odd good prime,
\[
 \frac{g(z)}{g(z^p)}=(1-z)^{(p-1)/2}\pmod p.
\]
Multiplication by the standard $F_{s_\ell}$ gives a polynomial
Frobenius quotient of degree $d_\ell(p)+(p-1)/2\leq p-1$.
Thus this nonrational algebraic gauge preserves Lucas congruences.
It does not preserve the first derivative at a regular evaluation point.
\end{remark}

Even preserving finite jets does not rule out algebraic gauges if
Lucas congruences are required only on a progression of primes.
The following construction states this limitation precisely.

\begin{proposition}\label{prop:C-algebraic-jet-lucas}
Let $M\geq2$, $N\geq1$, $D=M+N$, and let $m\geq2D$ be an integer.
Put
\[
 P(z)=1+2^{-M-1}z^N(4z-1)^M,\qquad
 g(z)=P(z)^{-1/m},\qquad H(z)=g(z)F_{1/2}(z),
\]
using the branch with $g(0)=1$.
The difference $H-F_{1/2}$ vanishes to order at least $N$ at zero
and at least $M$ at $z=1/4$. An ordinarily absolutely convergent
identity holds:
\[
 (1+6\thetaop)H(1/4)=\frac4\pi.
\]
For every sufficiently large prime $p\equiv1\pmod m$, $H$ satisfies
Lucas congruences modulo $p$.
\end{proposition}
\begin{proof}
On $|z|\leq1/4$ one has
\[
 |P(z)-1|\leq 2^{-M-1}4^{-N}2^M=2^{-1}4^{-N}<1.
\]
Hence the branch $g$ is analytic on a disk of radius greater than
$1/4$. The chosen logarithm gives $g(1/4)=1$, and the stated
vanishing orders follow from those of $P-1$.
In particular $g'(1/4)=0$, so the standard row $(A,B,C)=(1,6,256)$
is preserved.
Its series still converges absolutely there.
The gauge is nonrational: its pole-order equation $g^{-m}=P$ at
infinity would otherwise require $m$ to divide $D$, whereas $0<D<m$.

At a good prime $p\equiv1\pmod m$, the coefficients of $g$ are
$p$-integral: the binomial expansion has exponent $-1/m\in\Z_p$.
Reduction gives $g(z^p)=g(z)^p$, and therefore
\[
 \frac{H(z)}{H(z^p)}
 =T_p(z)P(z)^{(p-1)/m},\qquad
 T_p(z)=\sum_{n=0}^{p-1}c_{1/2}(n)z^n.
\]
This is a polynomial of degree
\[
 \frac{p-1}{2}+D\frac{p-1}{m}\leq p-1.
\]
Its functional identity with $H(z^p)$ is equivalent to Lucas
congruences. This proves the assertion.
\end{proof}

These constructions change the generating function and retain the CM
fiber at $z=1/4$; they are not counterexamples to standard CM forcing.
The exact factorial recurrence and its Lucas structure survive as
additional data that a successful transfer argument may need to use.
Neither their preservation nor the quotient covariance above supplies
the missing equality with an exact characteristic-zero Frobenius
splitting.

\subsection{Algebraic gauges on a majority of primes}
\label{subsec:C-majority-lucas-gauges}

The progression restriction in Proposition~\ref{prop:C-algebraic-jet-lucas}
cannot be removed. In fact, an upper prime density greater than one half
already classifies all algebraic gauges preserving Lucas congruences.
Here and below the density is taken relative to all rational primes.

\begin{theorem}[Classification of algebraic Lucas gauges]
\label{thm:C-majority-lucas-gauges}
Fix $s\in\{1/6,1/4,1/3,1/2\}$, and let
$g\in\Q[[z]]$ be algebraic over $\Q(z)$ with $g(0)=1$.
Let $\mathcal P_g$ be the set of primes at which the coefficients of
$gF_s$ are $p$-integral and their reductions satisfy Lucas congruences.
Then the following conditions are equivalent:
\begin{enumerate}
\item $\mathcal P_g$ has upper natural prime density greater than $1/2$;
\item $\mathcal P_g$ contains every sufficiently large prime;
\item $g(z)=(1-az)^{-1/2}$ for some $a\in\Q$, with the branch
normalized at zero.
\end{enumerate}
If these conditions hold and $g(z_0)=1$ at any nonzero regular point
on an analytic continuation of the branch, then $g=1$.
The condition $g'(0)=0$ likewise forces $g=1$.
\end{theorem}
\begin{proof}
Put $d=[\Q(z,g):\Q(z)]$. We first justify preservation of this degree
on reduction. Eisenstein's theorem on algebraic power series makes
all coefficients of $g$ integral outside finitely many primes.
The field $\Q(z,g)$ embeds in $\Q((z))$, whose elements algebraic
over $\Q$ are rational constants: differentiate a separable polynomial
equation to see that such a Laurent series has derivative zero.
Thus $\Q(z,g)/\Q$ is regular, and a primitive polynomial model
$H(z,Y)$ for $g$ is absolutely irreducible.
Absolute irreducibility persists outside finitely many primes.
Indeed, after homogenization, each possible factorization degree
is the image of multiplication between projective spaces of forms;
these images are closed because the maps are proper. The section
defined by $H$ avoids their union generically, hence meets it at
only finitely many primes. Excluding degree drops as well shows
\begin{equation}\label{eq:C-gauge-degree-reduction}
 [\mathbb F_p(z,g_p):\mathbb F_p(z)]=d
\end{equation}
for every sufficiently large $p$.

Suppose condition (1) holds. For a sufficiently large $p\in\mathcal P_g$,
write the Lucas identities for $f_p=(F_s)_p$ and $h_p=(gF_s)_p$ as
\[
 f_p(z)=T_p(z)f_p(z^p),\qquad
 h_p(z)=U_p(z)h_p(z^p),\qquad \deg T_p,\deg U_p<p.
\]
Their constant terms are one. Cancellation gives
\begin{equation}\label{eq:C-algebraic-gauge-kummer}
 g_p^{p-1}=T_p/U_p\in\mathbb F_p(z).
\end{equation}
All roots of $X^{p-1}-T_p/U_p$ lie in $\mathbb F_p(z,g_p)$,
since $\mu_{p-1}\subset\mathbb F_p$. This polynomial is separable,
so the field extension is Galois and its group embeds in
$\mathbb F_p^*$ through $\sigma\mapsto\sigma(g_p)/g_p$.
Equation~\eqref{eq:C-gauge-degree-reduction} therefore implies
$d\mid p-1$. Apart from finitely many exceptions,
$\mathcal P_g$ is contained in the progression $p\equiv1\pmod d$.
The prime number theorem in arithmetic progressions gives this
progression natural prime density $1/\varphi(d)$.
Since $\varphi(d)\geq2$ for $d>2$, condition (1) forces $d\leq2$.

If $d=1$, Proposition~\ref{prop:C-lucas-rational-rigidity} gives $g=1$.
If $d=2$, the nontrivial Galois automorphism at each sufficiently
large $p\in\mathcal P_g$ sends $g_p$ to $-g_p$.
The trace of $g$ over $\Q(z)$ consequently reduces to zero at
infinitely many primes and is zero. Hence
\[
                         g^2=R\in\Q(z).
\]

We need a uniform multiplicity bound for the roots of $T_p$.
By Proposition~\ref{prop:C-standard-lucas},
\[
 d_p:=\deg T_p=\left\lfloor\frac{p-1}{k_\ell}\right\rfloor,
 \qquad 1\leq d_p\leq(p-1)/2.
\]
The standard differential operator is
\[
 L_s=\thetaop^3-z(\thetaop+1/2)(\thetaop+s)(\thetaop+1-s).
\]
Since derivatives annihilate $f_p(z^p)$, the identity
$f_p=T_pf_p(z^p)$ and $L_sf_p=0$ imply $L_sT_p=0$.
Its leading ordinary-derivative coefficient is $z^3(1-z)$.
At any point other than $0,1$, a root of multiplicity $r\geq3$
would contribute a nonzero lowest term proportional to
$r(r-1)(r-2)$, because $r\leq d_p<p$. Thus its multiplicity
is at most two. At $z=1$, write $t=z-1$.
The coefficients of $\partial_z^3$ and $\partial_z^2$ are
$-t+O(t^2)$ and $-3/2+O(t)$, respectively.
A root of multiplicity $r\geq2$ would force
\[
                         r(r-1)(r-1/2)=0\pmod p.
\]
The only additional positive residue is $(p+1)/2>d_p$, so the
multiplicity at $1$ is at most one. Finally $T_p(0)=1$.

Write $R=P/Q$ in lowest terms, with $P(0)=Q(0)\ne0$.
At sufficiently large primes in $\mathcal P_g$,
\[
                       U_p=T_p(Q/P)^{(p-1)/2}.
\]
If $P$ had a root, polynomiality would require a root of $T_p$
of multiplicity at least $(p-1)/2>2$, impossible for $p\geq7$.
Thus $P$ is constant. Normalize to write $R=1/Q_0$ with $Q_0(0)=1$.
Then
\[
 \deg U_p=d_p+\tfrac{p-1}{2}\deg Q_0\leq p-1.
\]
Since $d_p>0$, one has $\deg Q_0\leq1$.
This gives $Q_0=1-az$ for $a\in\Q$ and proves (3).

Conversely, for such a gauge at each sufficiently large odd prime,
\[
 \frac{g(z)F_s(z)}{g(z^p)F_s(z^p)}
       =T_p(z)(1-az)^{(p-1)/2}\pmod p.
\]
This is a polynomial of degree at most $p-1$ and constant term one,
so it gives Lucas congruences. Thus (3) implies (2), which implies (1).
Finally $g(z_0)=1$ gives $1-az_0=1$, and $g'(0)=a/2$.
Both additional normalization assertions follow.
\end{proof}

\begin{remark}[Sharp density threshold and scope]
For any $a\in\Q^*$, the cubic gauge $g=(1-az)^{-1/3}$ has
Lucas-prime density exactly $1/2$. Its degree makes $p\equiv1\pmod3$
necessary at every sufficiently large Lucas prime. Conversely, at
each sufficiently large prime in that progression its Frobenius
quotient is
\[
                  T_p(z)(1-az)^{(p-1)/3},
\]
of degree at most $5(p-1)/6<p$, so Lucas holds.
This proves sharpness for the gauge classification, without an
additional target-value normalization. The theorem specifies the
standard section among its algebraic scalar gauges once that value
is normalized; it does not identify a complex period relation with
a Frobenius condition and therefore does not prove CM forcing.
\end{remark}

